\section*{Chapter \ref{ch:b3:glaciology} --- Glaciology}

\begin{solution}{exo:b3:glaciology:1}
$\tau_b = \rho gH\sin\alpha = 917\times9.81\times250\times\sin\ang{3} =
\qty{118}{kPa}$, about \qty{1.2}{bar}: the order of the stresses under real
glaciers.
\end{solution}

\begin{solution}{exo:b3:glaciology:2}
(a)~$\qty{2.9e15}{m^3}\times917 = \qty{2.66e18}{kg} = \qty{2.66e6}{Gt}$;
divided by \qty{360}{Gt/mm}, \qty{7390}{mm}, \qty{7.4}{m}, the value quoted.
(b)~$\qty{3.0e16}{m^3}\times917/360 = \qty{76}{m}$, against \qty{58}{m}.
Greenland's ice rests mostly on bedrock above sea level, so all of it would
add water to the ocean. Much of Antarctica's rests on a bed below sea
level: that ice already displaces sea water, and melting it only replaces
the water it displaces; the floating ice shelves add nothing at all.
\end{solution}

\begin{solution}{exo:b3:glaciology:3}
Gain $\num{6e6}\times1.2 = \qty{7.2e6}{m^3}$, loss $\num{4e6}\times2.1 =
\qty{8.4e6}{m^3}$ of water: net \qty{-1.2e6}{m^3}, a mean specific balance
of $-1.2\times10^6/10^7 = \qty{-0.12}{m}$ water equivalent. The
accumulation-area ratio is $6/10 = 0.6$. The glacier is losing mass:
shrinking.
\end{solution}

\begin{solution}{exo:b3:glaciology:4}
$1 - 820/917 = \num{0.106}$ and $1 - 830/917 = \num{0.095}$: about a tenth
of the volume is still pore space when the pores close.
\end{solution}

\begin{solution}{exo:b3:glaciology:5}
$\tau_b = 917\times9.81\times300\times\sin\ang{3} = \qty{141.2}{kPa}$. By
the slab theorem, $u_s = \tfrac{2A}{4}\tau_b^3H = 0.5\times\num{1.43e-16}
\times(\num{1.412e5})^3\times300 = \qty{60}{m/yr}$. Twice as thick: $u_s
\propto H^{n+1} = H^4$, sixteen times faster. Half the slope: $u_s \propto
(\sin\alpha)^3$, eight times slower.
\end{solution}

\begin{solution}{exo:b3:glaciology:6}
At \qty{-25}{\celsius} (\qty{248.15}{K}), $A = \num{3.61e-13}\exp(-60000/
(8.314\times248.15)) = \qty{8.5e-26}{Pa^{-3}.s^{-1}}$, or
\qty{2.7e-18}{Pa^{-3}.yr^{-1}}. At \qty{-10}{\celsius}, $\num{3.61e-13}
\exp(-27.42) = \qty{4.4e-25}{Pa^{-3}.s^{-1}} = \qty{1.4e-17}{Pa^{-3}.yr^{-1}}$.
At \qty{0}{\celsius} (warm branch) \qty{1.4e-16}{Pa^{-3}.yr^{-1}}. Under
the same stresses the speeds are in the ratio of the \omterm{def:b3:glaciology:glen}{rate factors}:
$\num{1.43e-16}/\num{2.67e-18} = 54$.
\end{solution}

\begin{solution}{exo:b3:glaciology:7}
$h = \sqrt{2\times10^5\times10^6/(917\times9.81)} = \qty{4.7}{km}$, close to
the \qty{4.9}{km} of Antarctica: the plastic model gets the thickness of an
ice sheet from its width alone.
\end{solution}

\begin{solution}{exo:b3:glaciology:8}
$\tau = 200/5 = \num{40}\text{ years}$. The length would go on adjusting for
several decades, a few response times, after the warming stopped: today's
glaciers have not yet caught up with today's climate.
\end{solution}

\begin{solution}{exo:b3:glaciology:9}
(a)~$h = \sqrt{2\times2.03\times20\times\num{8.64e6}/(917\times\num{3.34e5})}
= \qty{1.51}{m}$. Since $h^2 \propto t$, \qty{3}{m} takes $100\times
(3/1.51)^2 = 393$ days, more than a winter. (b)~Ice much thicker than
\qty{3}{m} cannot grow by freezing in one winter: it is ice that has
survived several, or ice piled up by the wind and the currents in ridges.
\end{solution}

\begin{solution}{exo:b3:glaciology:10}
(a)~$H = (1025/917)\times800 = \qty{894}{m}$. (b)~At \qty{900}{m} the
grounding-line thickness is \qty{1006}{m}, $9/8$ of the former; the flux
rises by $(9/8)^5 = 1.8$. (c)~As the line retreats down a bed deepening
inland, the ice at the line is thicker, discharges more, thins the ice
behind, and the line retreats further: a positive feedback. An ice shelf
buttressing the grounded ice reduces the flux at the \omterm{def:b3:glaciology:grounding-line}{grounding line} and can
hold it in place even on a downward slope.
\end{solution}

\begin{solution}{exo:b3:glaciology:11}
With $h(r) = c\sqrt{R - r}$, $c = \sqrt{2\tau_0/(\rho g)}$, $V =
\int_0^R2\pi rc\sqrt{R - r}\,\mathrm dr$. Substituting $s = R - r$,
\[
\int_0^R(R - s)\sqrt s\,\mathrm ds = \tfrac23R^{5/2} - \tfrac25R^{5/2} =
\tfrac4{15}R^{5/2},
\]
so $V = \frac{8\pi}{15}cR^{5/2}$. Doubling $R$
multiplies $V$, and the \omterm{def:b3:glaciology:sle}{sea-level equivalent}, by $2^{5/2} = 5.7$.
\end{solution}

\begin{solution}{exo:b3:glaciology:12}
(a)~March: $14.13/16.34 - 1 = \qty{-13.5}{\percent}$; September: $4.81/7.05
- 1 = \qty{-31.8}{\percent}$. (b)~In March the ice edge is held largely by
the coasts that surround the Arctic Ocean, and refreezing each winter
fills the basin again; the thinner ice of recent years melts out more
completely each summer. And the ice--albedo feedback acts in summer: open
water absorbs about \qty{94}{\percent} of the sunlight, ice reflects
\qtyrange{50}{70}{\percent}, so each loss of summer ice warms the water and
melts more.
\end{solution}

\begin{solution}{pb:b3:glaciology:1}
\textbf{1.} The column above the bed, of weight $\rho gH$ per unit area, is
pushed downslope by the component $\rho gH\sin\alpha$, balanced by the
shear stress on its base; for small slopes $\sin\alpha \approx \alpha$ (the
slab theorem with the surface slope).

\textbf{2.} $917\times9.81\times580\times0.0190 = \qty{99.1}{kPa}$;
$917\times9.81\times470\times0.0240 = \qty{101.5}{kPa}$; $917\times9.81
\times335\times0.0330 = \qty{99.4}{kPa}$.

\textbf{3.} The stress is the same within \qty{2}{\percent} while the
thickness almost halves: the cap is close to plastic, with $\tau_0 \approx
\qty{100}{kPa}$.

\textbf{4.} $2\times10^5/(917\times9.81) = \qty{22.23}{m}$.

\textbf{5.} $h = \sqrt{22.23\,(R - r)}$: \qty{577}{m}, \qty{472}{m} and
\qty{333}{m} at 15, 10 and \qty{5}{km} from the margin, within a few metres
of the soundings.

\textbf{6.} $\sqrt{22.23\times20000} = \qty{667}{m}$.

\textbf{7.} As in \cref{exo:b3:glaciology:11}: $V = \int_0^R2\pi r
\sqrt{22.23(R - r)}\,\mathrm dr = \frac{8\pi}{15}\sqrt{2\tau_0/(\rho
g)}\,R^{5/2}$.

\textbf{8.} $\frac{8\pi}{15}\times4.715\times(\num{2e4})^{5/2} =
\qty{4.47e11}{m^3} = \qty{447}{km^3}$.

\textbf{9.} $\pi\times20^2 = \qty{1257}{km^2}$; mean thickness
$447/1257 = \qty{0.356}{km}$, \qty{356}{m}, a little over half the
central thickness.

\textbf{10.} $2^{5/2} = 5.7$ times.

\textbf{11.} $A = \num{3.61e-13}\exp(-60000/(8.314\times263.15)) =
\qty{4.45e-25}{Pa^{-3}.s^{-1}}$, times \qty{3.156e7}{s}:
\qty{1.40e-17}{Pa^{-3}.yr^{-1}}.

\textbf{12.} At \qty{10}{km}, $H = \qty{471.5}{m}$, slope $471.5/20000 =
\num{0.0236}$, $\tau_b = \qty{100}{kPa}$; $u_s = \tfrac{2A}{4}\tau_b^3H =
0.5\times\num{1.40e-17}\times10^{15}\times471.5 = \qty{3.3}{m/yr}$.

\textbf{13.} $q = \tfrac{2A}{5}\tau_b^3H^2 = 0.4\times\num{1.40e-17}
\times10^{15}\times471.5^2 = \qty{1250}{m^2/yr}$; mean velocity $q/H =
\qty{2.6}{m/yr}$, four-fifths of the surface speed.

\textbf{14.} $2\pi rq = \pi r^2\bar b$, so $\bar b = 2q/r = 2\times1250/10^4
= \qty{0.25}{m}$ of ice a year over the inner \qty{10}{km}.

\textbf{15.} $h = 450$ where $R - r = 450^2/22.23 = \qty{9.11}{km}$, at
$r = \qty{10.9}{km}$; the accumulation area is the disc inside it, $(10.9/
20)^2 = 0.30$ of the cap.

\textbf{16.} $A$ rises by \num{10.2}, and so do the speeds under the same
stresses. The ice would flow out faster than the balance supplies it: the
cap would thin and spread until the stresses, lower, again matched the
balance; a softer plastic cap has a smaller yield stress and a flatter
profile.

\textbf{17.} $\num{4.47e11}\times917 = \qty{4.10e14}{kg} = \qty{410}{Gt}$.

\textbf{18.} $410/360 = \qty{1.14}{mm}$.

\textbf{19.} $356/0.5 = 711$ years.

\textbf{20.} As the surface lowers it lies in warmer air and melts faster,
the elevation--mass balance feedback; and as it thins the stresses and the
flow change too. The thinning accelerates.

\textbf{21.} $410/243 = \num{1.7}\text{ years}$: the whole cap is less than two
years of Greenland's loss in 2010--2019.

\textbf{22.} The Arkel ice cap holds a \omterm{def:b3:glaciology:sle}{sea-level equivalent} of
\textbf{\qty[text-series-to-math=true]{1.1}{mm}} (\qty{1.14}{mm}).
\end{solution}
